\chapter{Bridges, Articulation Points, and Biconnectivity}

A connected graph can still be fragile. Pull out one wire and the network splits;
cut one road and two towns can no longer reach each other. The vertices and edges
whose removal disconnects a graph are exactly the structural weak points, and a
single depth-first traversal finds all of them. The engine behind that traversal is
the \emph{low-link} value, a number Robert Tarjan attached to each vertex in 1972
that quietly encodes ``what is the highest ancestor I can reach without using the
tree edge from my parent?'' Once you internalize that one quantity, bridges,
articulation points, biconnected components, and the block-cut tree all fall out of
the same DFS in linear time.

This chapter builds the low-link machinery slowly, traces it by hand on a small
graph, proves the bridge and articulation-point characterizations, assembles the
biconnected components into a block-cut tree, and then steps back to the cleaner
combinatorial statements: $2$-edge-connectivity and $2$-vertex-connectivity, Menger's
theorem in its low-connectivity form, and what all of this says about network
reliability.

\chapterline

\section{What we are looking for}
\label{ch10:setup}

Throughout, $\Gph=(V,E)$ is an \emph{undirected} graph with $n=\abs{V}$ vertices and
$m=\abs{E}$ edges. We assume $\Gph$ is connected unless stated otherwise; if it is
not, every algorithm below runs once per connected component and we union the
results.

\begin{definition}[Articulation point]
\label{ch10:def-artic}
A vertex $v\in V$ is an \emph{articulation point} (or \emph{cut vertex}) if deleting
$v$ together with all edges incident to it increases the number of connected
components of $\Gph$. Equivalently, there exist vertices $x,y\neq v$ such that every
path from $x$ to $y$ passes through $v$.
\end{definition}

\begin{definition}[Bridge]
\label{ch10:def-bridge}
An edge $e\in E$ is a \emph{bridge} (or \emph{cut edge}) if deleting $e$ (but keeping
both endpoints) increases the number of connected components. Equivalently, $e=\{u,w\}$
is a bridge iff there is no path from $u$ to $w$ in $\Gph-e$, i.e.\ $e$ lies on no
cycle.
\end{definition}

\begin{intuitionbox}
A bridge is an edge with no ``backup route.'' If you can get from one endpoint to the
other by some other path, that path plus the edge forms a cycle, and removing one
edge of a cycle never disconnects anything. So \emph{bridges are precisely the edges
that lie on no cycle}. Articulation points are the vertex analogue: a vertex is a cut
vertex if some pair of its neighbors can only see each other ``through'' it.
\end{intuitionbox}

These two notions are siblings but not identical. The endpoint of a bridge is often,
but not always, an articulation point, and an articulation point need not be incident
to any bridge. The smallest illuminating example is a ``bowtie'': two triangles
sharing a single vertex.

\begin{center}
\begin{tikzpicture}[every node/.style={circle,draw,fill=blue!8,minimum size=7mm,inner sep=0pt,font=\small},thick]
  \node (a) at (0,0.9) {$a$};
  \node (b) at (0,-0.9) {$b$};
  \node[fill=trapred!25] (c) at (1.6,0) {$c$};
  \node (d) at (3.2,0.9) {$d$};
  \node (e) at (3.2,-0.9) {$e$};
  \draw (a)--(b); \draw (a)--(c); \draw (b)--(c);
  \draw (c)--(d); \draw (c)--(e); \draw (d)--(e);
\end{tikzpicture}
\end{center}

Here $c$ is an articulation point (deleting it separates $\{a,b\}$ from $\{d,e\}$), yet
\emph{no} edge is a bridge, because every edge lies on a triangle. Conversely, a single
edge joining two triangles is a bridge, and both of its endpoints are articulation
points. Keep this asymmetry in mind: the algorithms for the two problems share an
engine but differ in one comparison ($<$ versus $\le$).

\begin{definition}[Biconnected; $2$-connected]
\label{ch10:def-biconn}
A connected graph with at least three vertices is \emph{biconnected} (equivalently
\emph{$2$-vertex-connected}) if it has no articulation point: removing any single
vertex leaves it connected. A connected graph is \emph{$2$-edge-connected} if it has
no bridge: removing any single edge leaves it connected.
\end{definition}

\begin{focusbox}
Three connectivity questions, one DFS:
\begin{itemize}
  \item Which \emph{edges} are bridges? ($2$-edge-connectivity is ``no bridges.'')
  \item Which \emph{vertices} are articulation points? ($2$-vertex-connectivity is
        ``no cut vertices.'')
  \item How does the graph decompose into maximal pieces that have neither weakness
        (the \emph{biconnected components}), and how do those pieces fit together (the
        \emph{block-cut tree})?
\end{itemize}
All three are answered by computing, for every vertex, two integers: its discovery
time and its low-link value.
\end{focusbox}

\section{The DFS tree and why it has no cross edges}
\label{ch10:dfstree}

Everything rests on one structural fact about depth-first search in undirected graphs.
Run DFS from a root $r$. As we explore, we assign each vertex a \emph{discovery time}
$\mathrm{disc}[v]$, the order in which DFS first visits it ($1,2,3,\dots$). The edges
split into two kinds relative to the resulting rooted tree:

\begin{itemize}
  \item \textbf{Tree edges}: the edge $\{u,v\}$ along which DFS first reached $v$ from
        $u$. These form the DFS tree $T$.
  \item \textbf{Back edges}: every other edge. Such an edge $\{u,v\}$ always joins a
        vertex to one of its \emph{ancestors} in $T$.
\end{itemize}

\begin{whybox}
\textbf{No cross edges in undirected DFS.} Suppose, when DFS is exploring $u$, it looks
at a non-tree edge $\{u,v\}$. Because the edge exists, $v$ was discovered at some
point. If $v$ were neither an ancestor nor a descendant of $u$, then $v$ would have
been fully finished before $u$ started or vice versa; but if $v$ finished before $u$
started, the edge $\{u,v\}$ would have been examined from $v$'s side and would have
made $u$ a descendant of $v$ --- contradiction. Working through the two cases (which
endpoint is discovered first) shows the only possibility for a non-tree edge is that
one endpoint is an ancestor of the other. Hence every non-tree edge of an undirected
DFS is a \emph{back edge}. This is the property that makes low-link work; it fails for
directed graphs, where cross and forward edges appear.
\end{whybox}

Because there are no cross edges, the cycle structure of $\Gph$ is completely captured
by back edges: every cycle uses exactly one back edge together with the tree path
between its endpoints. This is the lever we pull.

\begin{center}
\begin{tikzpicture}[
  every node/.style={circle,draw,fill=blue!8,minimum size=8mm,inner sep=0pt,font=\small},
  treeedge/.style={very thick,nodeblue},
  backedge/.style={trapred,dashed,thick}]
  \node (r) at (0,0) {$r$};
  \node (a) at (-2,-1.4) {$a$};
  \node (b) at (2,-1.4) {$b$};
  \node (c) at (-2,-2.8) {$c$};
  \node (d) at (0,-2.8) {$d$};
  \node (e) at (2,-2.8) {$e$};
  \draw[treeedge] (r)--(a); \draw[treeedge] (r)--(b);
  \draw[treeedge] (a)--(c); \draw[treeedge] (a)--(d);
  \draw[treeedge] (b)--(e);
  \draw[backedge] (c) to[bend left=20] (r);
  \draw[backedge] (d) to[bend right=25] (r);
  \node[draw=none,fill=none,font=\scriptsize\itshape,nodeblue] at (-3.6,-0.7) {tree edge};
  \node[draw=none,fill=none,font=\scriptsize\itshape,trapred] at (3.4,-2.2) {back edge};
\end{tikzpicture}
\end{center}

The solid edges are the tree; the dashed edges go from a vertex up to a proper
ancestor. The back edge $c\!\to\!r$ creates the cycle $r,a,c,r$; the back edge
$d\!\to\!r$ creates the cycle $r,a,d,r$. Notice $b\!\to\!e$ is a tree edge with
\emph{no} back edge reaching above $b$: that tree edge is a bridge.

\section{Low-link: the highest reachable ancestor}
\label{ch10:lowlink}

\begin{definition}[Low-link]
\label{ch10:def-low}
Fix a DFS tree. For a vertex $v$, define
\[
  \mathrm{low}[v] \;=\; \min\Big(
    \{\mathrm{disc}[v]\}
    \;\cup\;
    \{\mathrm{disc}[w] : \{v,w\}\text{ a back edge}\}
    \;\cup\;
    \{\mathrm{low}[c] : c\text{ a child of }v\text{ in }T\}
  \Big).
\]
In words, $\mathrm{low}[v]$ is the smallest discovery time reachable from $v$ by going
down through zero or more tree edges in the subtree of $v$ and then taking at most one
back edge.
\end{definition}

\begin{intuitionbox}
Read $\mathrm{low}[v]$ as: ``Starting at $v$, staying inside my own subtree, and
allowed to climb exactly one back edge at the end, how high up the tree can I get?''
Low discovery time means high in the tree (the root has time $1$). If the highest I
can reach is strictly above my parent, my subtree is anchored to the rest of the graph
even without the edge from my parent --- there is a detour. If I cannot reach above my
parent, that single tree edge is the only thing tying my subtree in: cut it and I fall
off. That is the whole idea; the rest is bookkeeping.
\end{intuitionbox}

The definition is recursive and bottom-up, which is exactly how DFS finishes vertices:
a child's $\mathrm{low}$ is known by the time we return to the parent. We compute it
during the post-order phase of the same DFS that assigns discovery times.

\begin{algorithm}[H]
\caption{\textsc{DFS-LowLink}: discovery times and low-link in one pass}
\label{ch10:alg-lowlink}
\SetKwInOut{Input}{input}\SetKwInOut{Output}{output}
\Input{undirected graph $\Gph=(V,E)$ via adjacency lists; global counter $t\gets 0$}
\Output{arrays $\mathrm{disc}[\cdot]$ and $\mathrm{low}[\cdot]$}
\SetKwProg{Fn}{function}{}{}
\Fn{\textsc{Dfs}($u$, $p$)}{
  $t\gets t+1$;\quad $\mathrm{disc}[u]\gets t$;\quad $\mathrm{low}[u]\gets t$\;
  \ForEach{$w$ in $\mathrm{adj}[u]$}{
    \If{edge to $w$ is the (single) parent edge to $p$}{
      \textbf{continue}\tcp*{skip exactly one copy of the parent edge}
    }
    \eIf{$\mathrm{disc}[w]=0$}{
      \textsc{Dfs}($w$, $u$)\tcp*{tree edge}
      $\mathrm{low}[u]\gets\min(\mathrm{low}[u],\mathrm{low}[w])$\;
    }{
      $\mathrm{low}[u]\gets\min(\mathrm{low}[u],\mathrm{disc}[w])$\tcp*{back edge to ancestor $w$}
    }
  }
}
\ForEach{$v\in V$ with $\mathrm{disc}[v]=0$}{\textsc{Dfs}($v,\textsc{nil}$)}
\end{algorithm}

\begin{trapbox}
\textbf{The parent edge, and the multigraph subtlety.} When scanning $u$'s neighbors
we must not treat the edge back to the parent $p$ as a back edge --- otherwise
$\mathrm{low}[u]$ would always reach $\mathrm{disc}[p]$ and we would miss every bridge.
The naive guard ``if $w=p$ skip'' is \emph{wrong} when the graph has parallel edges: if
there are \emph{two} edges between $u$ and $p$, one is the tree edge and the other is a
genuine back edge that makes $\{u,p\}$ a non-bridge. The correct rule is to skip
exactly \emph{one} occurrence of the parent edge, which is why robust implementations
track edges by their index rather than by the neighbor label. With distinct neighbor
labels and a simple graph, the ``skip if $w=p$'' shortcut is fine; with multigraphs,
identify edges, not vertices.
\end{trapbox}

\begin{whybox}
\textbf{Correctness of the recurrence.} We prove by induction on finishing order that
when \textsc{Dfs}($u,\cdot$) returns, $\mathrm{low}[u]$ equals the minimum discovery
time over all vertices reachable from $u$ by a path that stays in $u$'s subtree except
for a single final back edge. The initialization sets $\mathrm{low}[u]=\mathrm{disc}[u]$
(the empty path). Each back edge $\{u,w\}$ to an ancestor contributes $\mathrm{disc}[w]$
directly. Each tree child $c$ is finished before $u$ returns, so by the inductive
hypothesis $\mathrm{low}[c]$ is already the correct subtree minimum for $c$; folding it
in via $\min$ extends the reachable set through $c$'s subtree. Since every in-subtree
``descend then climb once'' path either is empty, ends with a back edge from $u$
itself, or passes through exactly one child $c$ and continues inside $c$'s subtree, the
three cases are exhaustive. Hence the returned value is exactly the claimed minimum.
\end{whybox}

\subsection{Tracing low-link by hand}
\label{ch10:trace}

Let us run Algorithm~\ref{ch10:alg-lowlink} on a concrete graph and watch the numbers.
Take vertices $\{1,\dots,8\}$ with edges
\[
  \{1,2\},\{1,3\},\{2,3\},\{3,4\},\{4,5\},\{4,6\},\{5,6\},\{6,7\},\{7,8\}.
\]

\begin{center}
\begin{tikzpicture}[every node/.style={circle,draw,fill=blue!8,minimum size=8mm,inner sep=0pt,font=\small},thick]
  \node (1) at (0,0) {$1$};
  \node (2) at (-1.2,-1.3) {$2$};
  \node (3) at (1.2,-1.3) {$3$};
  \node (4) at (1.2,-2.8) {$4$};
  \node (5) at (0,-4.1) {$5$};
  \node (6) at (2.4,-4.1) {$6$};
  \node (7) at (3.6,-5.2) {$7$};
  \node (8) at (4.8,-6.1) {$8$};
  \draw (1)--(2); \draw (1)--(3); \draw (2)--(3);
  \draw[trapred,very thick] (3)--(4);
  \draw (4)--(5); \draw (4)--(6); \draw (5)--(6);
  \draw[trapred,very thick] (6)--(7); \draw[trapred,very thick] (7)--(8);
\end{tikzpicture}
\end{center}

The graph is a triangle $\{1,2,3\}$, joined by edge $\{3,4\}$ to a triangle
$\{4,5,6\}$, with a tail $6$--$7$--$8$ hanging off. The thick red edges
$\{3,4\}$, $\{6,7\}$, $\{7,8\}$ are the bridges, as we will confirm. We run DFS from
vertex $1$, always visiting neighbors in increasing label order. The tree edges
discovered are $1\!\to\!2$, $2\!\to\!3$, $3\!\to\!4$, $4\!\to\!5$, $5\!\to\!6$,
$6\!\to\!7$, $7\!\to\!8$; the back edges encountered are $3\!\to\!1$ (closing the first
triangle) and $6\!\to\!4$ (closing the second).

\begin{center}
\begin{tabular}{c|cccccccc}
\toprule
$v$            & $1$ & $2$ & $3$ & $4$ & $5$ & $6$ & $7$ & $8$\\
\midrule
$\mathrm{disc}[v]$ & $1$ & $2$ & $3$ & $4$ & $5$ & $6$ & $7$ & $8$\\
$\mathrm{low}[v]$  & $1$ & $1$ & $1$ & $4$ & $4$ & $4$ & $7$ & $8$\\
\bottomrule
\end{tabular}
\end{center}

Walk the post-order. Vertex $8$ is a leaf with no back edge: $\mathrm{low}[8]=8$. Then
$7$: its only non-parent neighbor is $8$ (a child), so $\mathrm{low}[7]=\min(7,
\mathrm{low}[8])=7$. Then $6$: it has child $7$ ($\mathrm{low}[7]=7$) and a back edge to
$4$ ($\mathrm{disc}[4]=4$), so $\mathrm{low}[6]=\min(6,7,4)=4$. Vertex $5$: child $6$ gives
$\mathrm{low}[5]=\min(5,4)=4$. Vertex $4$: child $5$ gives $\mathrm{low}[4]=\min(4,4)=4$.
Vertex $3$: child $4$ ($\mathrm{low}[4]=4$) and back edge to $1$ ($\mathrm{disc}[1]=1$),
so $\mathrm{low}[3]=\min(3,4,1)=1$. Vertex $2$: child $3$ gives $\min(2,1)=1$. Vertex
$1$: child $2$ gives $\min(1,1)=1$.

\begin{labbox}
Redo this trace but start DFS from vertex $4$ and visit neighbors in \emph{decreasing}
label order. The discovery times change completely, the tree reshapes, yet the set of
bridges you read off ($\{3,4\},\{6,7\},\{7,8\}$) and the set of articulation points
must come out identical. Bridges and cut vertices are properties of $\Gph$, not of the
traversal --- a good sanity check that you have implemented the comparisons correctly.
\end{labbox}

\section{Finding bridges}
\label{ch10:bridges}

\begin{theorem}[Bridge criterion]
\label{ch10:thm-bridge}
Let $\{u,v\}$ be a \emph{tree edge} with $v$ the child. Then $\{u,v\}$ is a bridge if
and only if
\[
  \mathrm{low}[v] > \mathrm{disc}[u].
\]
No back edge is ever a bridge.
\end{theorem}

\begin{whybox}
\textbf{Why the comparison is strict.} The subtree rooted at $v$ is attached to the
rest of $\Gph$ only through the tree edge $\{u,v\}$ \emph{unless} some vertex in $v$'s
subtree has a back edge climbing to $u$ or above $u$. The quantity $\mathrm{low}[v]$ is
exactly the highest discovery time any such climb reaches. ``Above $u$ or to $u$''
means $\mathrm{low}[v]\le\mathrm{disc}[u]$ --- a detour exists, no bridge. ``Strictly
below $u$,'' i.e.\ $\mathrm{low}[v]>\mathrm{disc}[u]$, means every escape from the
subtree must use the edge $\{u,v\}$ itself, so deleting it severs the subtree. A back
edge $\{x,y\}$ can never be a bridge because it closes a cycle (the tree path from $x$
to $y$ plus the edge), and no edge of a cycle is a bridge.
\end{whybox}

\begin{formalbox}
\textbf{Proof of Theorem~\ref{ch10:thm-bridge}.}
($\Leftarrow$) Suppose $\mathrm{low}[v]>\mathrm{disc}[u]$. Let $S$ be the set of
vertices in $v$'s subtree. Every edge of $\Gph$ with one endpoint in $S$ is either the
tree edge $\{u,v\}$, an internal tree edge of the subtree, or a back edge from some
$x\in S$ to an ancestor $y$. For such a back edge, $\mathrm{disc}[y]\ge\mathrm{low}[x]
\ge\mathrm{low}[v]>\mathrm{disc}[u]$, so $y$ is a \emph{descendant} of $u$ strictly below
$u$; since $y$ is an ancestor of $x\in S$ and lies below $u$, in fact $y\in S$.
Therefore the only edge crossing from $S$ to $V\setminus S$ is $\{u,v\}$. Removing it
disconnects $S$ from $u$; it is a bridge.

($\Rightarrow$) Suppose $\mathrm{low}[v]\le\mathrm{disc}[u]$. By definition of
$\mathrm{low}$ there is a vertex $x$ in $v$'s subtree with a back edge to an ancestor
$y$ satisfying $\mathrm{disc}[y]=\mathrm{low}[v]\le\mathrm{disc}[u]$, so $y$ is $u$ or a
proper ancestor of $u$. Then the path
\[
  u \;\to\; v \;\rightsquigarrow\; x \;\to\; y \;\rightsquigarrow\; u
\]
(tree edge, tree path down to $x$, back edge to $y$, tree path from $y$ back to $u$)
is a closed walk containing $\{u,v\}$ and using it once. Extracting a simple cycle
through $\{u,v\}$ shows the edge lies on a cycle, hence is not a bridge. $\qed$
\end{formalbox}

\begin{algorithm}[H]
\caption{\textsc{Bridges}: list all bridges in one DFS}
\label{ch10:alg-bridges}
\SetKwInOut{Output}{output}
\Output{set $B$ of bridge edges}
$B\gets\emptyset$;\quad $t\gets0$;\quad $\mathrm{disc}[\cdot]\gets0$\;
\SetKwProg{Fn}{function}{}{}
\Fn{\textsc{Dfs}($u$, $p_e$) \tcp*{$p_e$ = edge index used to enter $u$}}{
  $t\gets t+1$;\ $\mathrm{disc}[u]\gets\mathrm{low}[u]\gets t$\;
  \ForEach{edge $e=\{u,w\}$ in $\mathrm{adj}[u]$ with index $i$}{
    \lIf{$i=p_e$}{\textbf{continue}}
    \eIf{$\mathrm{disc}[w]=0$}{
      \textsc{Dfs}($w$, $i$); \ $\mathrm{low}[u]\gets\min(\mathrm{low}[u],\mathrm{low}[w])$\;
      \lIf{$\mathrm{low}[w]>\mathrm{disc}[u]$}{$B\gets B\cup\{e\}$}
    }{
      $\mathrm{low}[u]\gets\min(\mathrm{low}[u],\mathrm{disc}[w])$\;
    }
  }
}
\lForEach{$v\in V$ with $\mathrm{disc}[v]=0$}{\textsc{Dfs}($v,-1$)}
\end{algorithm}

On the traced graph: edge $3\!\to\!4$ has $\mathrm{low}[4]=4>\mathrm{disc}[3]=3$
--- bridge. Edge $6\!\to\!7$ has $\mathrm{low}[7]=7>\mathrm{disc}[6]=6$ --- bridge.
Edge $7\!\to\!8$ has $\mathrm{low}[8]=8>\mathrm{disc}[7]=7$ --- bridge. Every other
tree edge fails the test (e.g.\ $5\!\to\!6$ has $\mathrm{low}[6]=4\le5$), matching the
red edges in the picture.

\section{Finding articulation points}
\label{ch10:artic}

The vertex version uses the same $\mathrm{low}$ array with two changes: a non-strict
comparison, and a special rule for the root.

\begin{theorem}[Articulation-point criterion]
\label{ch10:thm-artic}
Run DFS from a chosen root $r$.
\begin{enumerate}
  \item The root $r$ is an articulation point if and only if it has \textbf{two or more}
        children in the DFS tree.
  \item A non-root vertex $u$ is an articulation point if and only if it has a child
        $v$ in the DFS tree with
        \[
          \mathrm{low}[v] \;\ge\; \mathrm{disc}[u].
        \]
\end{enumerate}
\end{theorem}

\begin{intuitionbox}
Compare with the bridge test $\mathrm{low}[v]>\mathrm{disc}[u]$. For an edge to be a
bridge, the child's escape must land \emph{strictly below} $u$. For a \emph{vertex}
$u$ to be a cut vertex, the child's escape only needs to fail to climb \emph{above}
$u$ --- landing exactly \emph{at} $u$ ($\mathrm{low}[v]=\mathrm{disc}[u]$) still traps
the subtree behind $u$, because the only way out passes through $u$ itself. Hence
$\ge$ for vertices, $>$ for edges. The root is special because it has no parent edge to
escape ``upward'' through; it is a cut vertex exactly when it branches the tree into
two or more subtrees that can only meet at $r$.
\end{intuitionbox}

\begin{formalbox}
\textbf{Proof of Theorem~\ref{ch10:thm-artic}.}
\emph{Root.} If $r$ has children $v_1,v_2$ in $T$, no edge joins the subtree of $v_1$ to
that of $v_2$ (such an edge would be a non-tree, non-back edge --- a cross edge ---
which undirected DFS forbids). So every $v_1$--$v_2$ path uses $r$; deleting $r$
separates them. Conversely if $r$ has a single child, its whole subtree is the rest of
the graph and is connected without $r$.

\emph{Non-root $u$.} ($\Leftarrow$) Suppose a child $v$ has
$\mathrm{low}[v]\ge\mathrm{disc}[u]$. Every back edge leaving $v$'s subtree reaches a
vertex with discovery time $\ge\mathrm{disc}[u]$, i.e.\ $u$ or a descendant of $u$;
none reaches a proper ancestor of $u$. Thus every path from $v$'s subtree to the part
of $\Gph$ above $u$ passes through $u$. Deleting $u$ separates $v$'s subtree (nonempty,
not containing $u$) from $u$'s parent --- so $u$ is a cut vertex.
($\Rightarrow$) Suppose every child $v$ has $\mathrm{low}[v]<\mathrm{disc}[u]$. Then each
child's subtree has a back edge to a proper ancestor of $u$, giving a detour around $u$
that connects that subtree to the part above $u$. Removing $u$, all of $u$'s child
subtrees stay connected to the ancestor portion (which is connected among itself), so
the graph minus $u$ stays connected; $u$ is not a cut vertex. $\qed$
\end{formalbox}

\begin{algorithm}[H]
\caption{\textsc{Articulation}: list all cut vertices in one DFS}
\label{ch10:alg-artic}
\SetKwInOut{Output}{output}
\Output{set $A$ of articulation points}
$A\gets\emptyset$;\ $t\gets0$;\ $\mathrm{disc}[\cdot]\gets0$\;
\SetKwProg{Fn}{function}{}{}
\Fn{\textsc{Dfs}($u$, $p$, \textnormal{isRoot})}{
  $t\gets t+1$;\ $\mathrm{disc}[u]\gets\mathrm{low}[u]\gets t$;\ $\mathrm{kids}\gets0$\;
  \ForEach{$w\in\mathrm{adj}[u]$}{
    \lIf{$w=p$}{\textbf{continue}}
    \eIf{$\mathrm{disc}[w]=0$}{
      $\mathrm{kids}\gets\mathrm{kids}+1$;\ \textsc{Dfs}($w$, $u$, \textsc{false})\;
      $\mathrm{low}[u]\gets\min(\mathrm{low}[u],\mathrm{low}[w])$\;
      \lIf{\textnormal{not} isRoot \textnormal{and} $\mathrm{low}[w]\ge\mathrm{disc}[u]$}{$A\gets A\cup\{u\}$}
    }{
      $\mathrm{low}[u]\gets\min(\mathrm{low}[u],\mathrm{disc}[w])$\;
    }
  }
  \lIf{isRoot \textnormal{and} $\mathrm{kids}\ge2$}{$A\gets A\cup\{u\}$}
}
\lForEach{$v\in V$ with $\mathrm{disc}[v]=0$}{\textsc{Dfs}($v,\textsc{nil},\textsc{true}$)}
\end{algorithm}

On the traced graph, root $=1$ has a single child, so it is not a cut vertex. Checking
non-roots: $u=3$ has child $4$ with $\mathrm{low}[4]=4\ge\mathrm{disc}[3]=3$ ---
articulation point. $u=4$ has child $5$ with $\mathrm{low}[5]=4\ge\mathrm{disc}[4]=4$
--- articulation point. $u=6$ has child $7$ with $\mathrm{low}[7]=7\ge6$ --- articulation
point. $u=7$ has child $8$ with $\mathrm{low}[8]=8\ge7$ --- articulation point. The cut
vertices are $\{3,4,6,7\}$. Vertices $2$ and $5$ are not (their children climb above
them); $1$ and $8$ are not (root with one child; leaf).

\begin{trapbox}
\textbf{Do not apply the $\ge$ rule to the root.} If you forget the root special case
and use $\mathrm{low}[v]\ge\mathrm{disc}[r]$ on the root, every root with at least one
child reports as a cut vertex, since $\mathrm{disc}[r]=1$ is the global minimum and any
$\mathrm{low}[v]\ge1$. The root genuinely needs the ``two children'' test.
Symmetrically, do not count a child twice in a multigraph: increment $\mathrm{kids}$
only on actual tree edges.
\end{trapbox}

\begin{historybox}
The discovery-time/low-link DFS is due to \textbf{Robert Tarjan}, ``Depth-first search
and linear graph algorithms'' (1972), the same paper that gave linear-time strongly
connected components. Hopcroft and Tarjan packaged the biconnectivity routine as
``Algorithm 447'' (Communications of the ACM, 1973). Before this, testing
$2$-connectivity was done
by repeated reachability checks costing $\Oh(nm)$ or worse. Tarjan's insight --- that a
single number per vertex, computed in post-order, certifies the entire cycle structure
--- collapsed a whole family of problems to one linear-time traversal. The low-link
idea reappears verbatim in Tarjan's SCC algorithm and in Gabow's later refinements; it
is one of the most reused tricks in all of graph algorithms.
\end{historybox}

\section{Biconnected components and the block-cut tree}
\label{ch10:bcc}

Articulation points and bridges are the weak \emph{spots}; the biconnected components
are the maximal \emph{strong pieces} between them.

\begin{definition}[Biconnected component / block]
\label{ch10:def-block}
Define a relation on edges: $e\sim f$ if $e=f$ or $e$ and $f$ lie on a common simple
cycle. This is an equivalence relation, and its classes are the \emph{biconnected
components} (or \emph{blocks}). Each block is either a maximal $2$-connected subgraph
or a single bridge edge.
\end{definition}

\begin{whybox}
\textbf{Why ``share a cycle'' is an equivalence relation.} Reflexivity and symmetry are
immediate. Transitivity is the content: if $e,f$ share a cycle $C_1$ and $f,g$ share a
cycle $C_2$, then $e$ and $g$ share a cycle. This is a small but real lemma --- the two
cycles overlap in at least the edge $f$, and one can splice them along their shared
portion into a single cycle containing both $e$ and $g$. (Splicing along a shared path,
not just a shared edge, is what makes it work even when $C_1,C_2$ meet in more than
$f$.) Consequently blocks \emph{partition the edges}. Vertices, by contrast, are
\emph{not} partitioned: an articulation point belongs to several blocks at once --- it
is exactly the vertex shared by two or more of them.
\end{whybox}

So edges split cleanly into blocks, and the articulation points are precisely the
vertices lying in more than one block. A bridge is a block all by itself (a $K_2$ that
is on no cycle). This is the cleanest way to see the relationship: \emph{bridges are the
single-edge blocks; articulation points are the shared vertices between blocks}.

\subsection{Extracting blocks with an edge stack}
\label{ch10:bccstack}

The same DFS that computes $\mathrm{low}$ peels off blocks using an auxiliary stack of
edges. Push each edge when first traversed. When the articulation test
$\mathrm{low}[w]\ge\mathrm{disc}[u]$ fires at $u$ for child $w$, every edge currently on
the stack down to and including $\{u,w\}$ forms one complete block; pop them off.

\begin{algorithm}[H]
\caption{\textsc{BiconnectedComponents}: peel blocks with an edge stack}
\label{ch10:alg-bcc}
\SetKwInOut{Output}{output}
\Output{list of blocks, each a set of edges}
$t\gets0$;\ $\mathrm{disc}[\cdot]\gets0$;\ stack $S\gets$ empty\;
\SetKwProg{Fn}{function}{}{}
\Fn{\textsc{Dfs}($u$, $p$)}{
  $t\gets t+1$;\ $\mathrm{disc}[u]\gets\mathrm{low}[u]\gets t$\;
  \ForEach{$w\in\mathrm{adj}[u]$ via edge $e$}{
    \lIf{$e$ is the parent edge}{\textbf{continue}}
    \If{$\mathrm{disc}[w]=0$}{
      push $e$ onto $S$\;
      \textsc{Dfs}($w$, $u$);\ $\mathrm{low}[u]\gets\min(\mathrm{low}[u],\mathrm{low}[w])$\;
      \If{$\mathrm{low}[w]\ge\mathrm{disc}[u]$}{
        start a new block; pop edges from $S$ until $e$ is popped (inclusive); emit the block\;
      }
    }
    \ElseIf{$\mathrm{disc}[w]<\mathrm{disc}[u]$}{
      push $e$ onto $S$;\ $\mathrm{low}[u]\gets\min(\mathrm{low}[u],\mathrm{disc}[w])$\tcp*{back edge}
    }
  }
}
\lForEach{$v\in V$ with $\mathrm{disc}[v]=0$}{\textsc{Dfs}($v,\textsc{nil}$)}
\end{algorithm}

\begin{whybox}
\textbf{Why the stack respects block boundaries.} Edges are pushed in DFS order, so the
stack always holds a contiguous ``frontier'' of edges descending from the current
search path. When $\mathrm{low}[w]\ge\mathrm{disc}[u]$, the subtree of $w$ cannot reach
above $u$, so no future edge will ever connect into $w$'s subtree from higher up: every
edge of the block that $\{u,w\}$ belongs to has already been pushed and sits at the top
of the stack, ending exactly at $\{u,w\}$. Popping down to $\{u,w\}$ therefore emits one
complete block and nothing more. The condition is the \emph{same} $\ge$ test as for
articulation points, which is no accident: a block boundary at $u$ is precisely what
makes $u$ a cut vertex.
\end{whybox}

On our running graph the blocks come out as
\[
  \{12,13,23\}\ (\text{triangle}),\quad
  \{34\}\ (\text{bridge}),\quad
  \{45,46,56\}\ (\text{triangle}),\quad
  \{67\}\ (\text{bridge}),\quad
  \{78\}\ (\text{bridge}),
\]
five blocks: two triangles and three single-edge bridges. The cut vertices $3,4,6,7$
are exactly the vertices appearing in two blocks each.

\subsection{The block-cut tree}
\label{ch10:bctree}

The blocks and cut vertices assemble into a tree that records the coarse shape of the
graph.

\begin{definition}[Block-cut tree]
\label{ch10:def-bctree}
The \emph{block-cut tree} $\mathcal{T}$ of a connected graph $\Gph$ has one node per
block and one node per articulation point. Connect a cut-vertex node $c$ to a block
node $\beta$ iff $c$ is a vertex of $\beta$. The result is a tree (a forest if $\Gph$
is disconnected).
\end{definition}

\begin{whybox}
\textbf{Why it is acyclic.} Suppose $\mathcal{T}$ contained a cycle. That cycle
alternates between cut-vertex nodes and block nodes. Splicing the cycles inside the
participating blocks together along the shared cut vertices would merge them all into a
single biconnected piece --- contradicting that they were \emph{distinct maximal}
blocks. So $\mathcal{T}$ has no cycle, and since $\Gph$ is connected $\mathcal{T}$ is
connected; hence it is a tree. Its leaves are always blocks (a cut vertex has degree
$\ge2$ in $\mathcal{T}$ by definition), a fact used in many reliability and augmentation
arguments.
\end{whybox}

\begin{center}
\begin{tikzpicture}[
  blk/.style={rectangle,rounded corners,draw,fill=green!12,minimum size=7mm,inner sep=3pt,font=\small},
  cut/.style={circle,draw,fill=trapred!20,minimum size=7mm,inner sep=0pt,font=\small},
  thick]
  \node[blk] (B1) at (0,0) {$B_1{=}\{1,2,3\}$};
  \node[cut] (c3) at (2.6,0) {$3$};
  \node[blk] (B2) at (4.0,0) {$\{3,4\}$};
  \node[cut] (c4) at (5.4,0) {$4$};
  \node[blk] (B3) at (7.7,0.0) {$B_3{=}\{4,5,6\}$};
  \node[cut] (c6) at (7.7,-1.5) {$6$};
  \node[blk] (B4) at (7.7,-3.0) {$\{6,7\}$};
  \node[cut] (c7) at (7.7,-4.5) {$7$};
  \node[blk] (B5) at (7.7,-6.0) {$\{7,8\}$};
  \draw (B1)--(c3); \draw (c3)--(B2); \draw (B2)--(c4); \draw (c4)--(B3);
  \draw (B3)--(c6); \draw (c6)--(B4); \draw (B4)--(c7); \draw (c7)--(B5);
\end{tikzpicture}
\end{center}

The block-cut tree of the running graph is a path: triangle $B_1$ -- cut $3$ -- bridge
-- cut $4$ -- triangle $B_3$ -- cut $6$ -- bridge -- cut $7$ -- bridge $\{7,8\}$. Square
nodes are blocks, round red nodes are cut vertices. Reading the tree, the bridges are
exactly the single-edge block nodes, and the cut vertices are exactly the round nodes
--- the whole connectivity skeleton at a glance.

\begin{applicationbox}
The block-cut tree turns ``what fails if vertex $v$ dies?'' into a tree query. The
number of components created by deleting a cut vertex $c$ equals its degree in
$\mathcal{T}$. Two query vertices $x,y$ are separated by deleting $c$ iff $c$ lies on
the $\mathcal{T}$-path between the blocks containing $x$ and $y$. Many ``is there a
path avoiding $v$'' or ``which articulation point disconnects $x$ from $y$'' questions
reduce to LCA and path queries on $\mathcal{T}$ (see the LCA chapter), answering each in
$\Oh(\log n)$ or $\Oh(1)$ after linear preprocessing.
\end{applicationbox}

\section{Complexity}
\label{ch10:complexity}

\begin{theorem}
\label{ch10:thm-time}
Bridges, articulation points, all biconnected components, and the block-cut tree can
each be computed in $\Oh(n+m)$ time and $\Oh(n)$ extra space (plus $\Oh(m)$ for the
block lists or the edge stack).
\end{theorem}

\begin{formalbox}
\textbf{Proof.} All four algorithms are a single DFS over $\Gph$ with $\Oh(1)$ extra
work per edge examination. Each vertex is entered once (the $\mathrm{disc}$ guard) and
each edge is scanned at most twice, once from each endpoint, doing constant work:
updating a $\min$, comparing two integers, or pushing/popping the edge stack. Summing,
the traversal is $\Oh(n+m)$. The arrays $\mathrm{disc}$ and $\mathrm{low}$ use $\Oh(n)$
words; the edge stack holds at most $m$ edges and each edge is pushed and popped exactly
once, so its total handling cost is $\Oh(m)$ and it never exceeds $\Oh(m)$ space. The
block-cut tree has at most $n$ cut-vertex nodes and at most $m$ block nodes, so it is
built in $\Oh(n+m)$. $\qed$
\end{formalbox}

\begin{complexitybox}
\textbf{Cost summary.}
\begin{center}
\begin{tabular}{lll}
\toprule
Task & Time & Space (beyond input)\\
\midrule
Discovery times \& low-link & $\Oh(n+m)$ & $\Oh(n)$\\
All bridges & $\Oh(n+m)$ & $\Oh(n)$\\
All articulation points & $\Oh(n+m)$ & $\Oh(n)$\\
Biconnected components (edge stack) & $\Oh(n+m)$ & $\Oh(m)$\\
Block-cut tree & $\Oh(n+m)$ & $\Oh(n+m)$\\
\bottomrule
\end{tabular}
\end{center}
A non-recursive (explicit-stack) DFS is advisable in practice: on a path of $10^6$
vertices the recursive version overflows the call stack, while the connectivity facts
are unchanged.
\end{complexitybox}

\section{$2$-edge- and $2$-vertex-connectivity}
\label{ch10:2conn}

The local tests assemble into clean global statements. Recall a graph is
$2$-edge-connected if it has no bridge, and $2$-vertex-connected (biconnected, $n\ge3$)
if it has no articulation point.

\begin{corollary}
\label{ch10:cor-2ec}
A connected graph $\Gph$ with $m\ge1$ is $2$-edge-connected iff the bridge DFS reports
no bridge, iff every edge lies on a cycle, iff $\Gph$ has a single biconnected
component containing all edges. All of these are decidable in $\Oh(n+m)$.
\end{corollary}

\begin{corollary}
\label{ch10:cor-2vc}
A connected graph $\Gph$ with $n\ge3$ is $2$-vertex-connected iff the articulation DFS
reports no cut vertex, iff between every pair of vertices there are two
internally-vertex-disjoint paths, iff the block-cut tree is a single block node.
\end{corollary}

The ``two disjoint paths'' phrasing is the low-connectivity case of a deep theorem.

\begin{theorem}[Menger, edge and vertex forms]
\label{ch10:thm-menger}
For distinct non-adjacent $s,t$, the maximum number of pairwise edge-disjoint $s$--$t$
paths equals the minimum number of edges whose removal separates $s$ from $t$; and the
maximum number of internally-vertex-disjoint $s$--$t$ paths equals the minimum number
of vertices (other than $s,t$) whose removal separates them.
\end{theorem}

\begin{whybox}
\textbf{The $k=2$ special case is exactly our two tests.} Menger says local
connectivity equals a global path count. For $k=2$: ``no single edge separates any
pair'' (no bridge, $2$-edge-connected) is equivalent to ``every pair has $2$
edge-disjoint paths''; ``no single vertex separates any pair'' (no cut vertex,
$2$-vertex-connected) is equivalent to ``every pair has $2$ internally-disjoint
paths,'' which for a graph on $\ge3$ vertices means every pair lies on a common cycle.
The general theorem is proved via max-flow min-cut in the flow chapters; here the DFS
gives the $k=2$ case directly and in linear time, without any flow computation.
\end{whybox}

\begin{formalbox}
\textbf{Cycle characterization of biconnectivity.} We sketch why ``no articulation
point'' (with $n\ge3$, connected) is equivalent to ``every two vertices lie on a common
cycle.'' If every pair lies on a cycle, then for any $v$ and any pair $x,y\neq v$, take
a cycle through $x,y$; it gives an $x$--$y$ path avoiding $v$, so $v$ is not a cut
vertex. Conversely, suppose $\Gph$ is biconnected and take any $x,y$. Two
internally-disjoint $x$--$y$ paths (which exist by the vertex Menger bound, since no
single vertex separates them) concatenate into a cycle through $x$ and $y$. An
induction building the graph up by ``ear decompositions'' --- repeatedly adding a path
(an ``ear'') whose endpoints are already present --- makes this constructive: a graph
is $2$-connected iff it has an open ear decomposition starting from a cycle (Whitney's
theorem). $\qed$
\end{formalbox}

\begin{historybox}
\textbf{Hassler Whitney} (1932) introduced higher connectivity and proved the
ear-decomposition characterization of $2$-connected graphs. \textbf{Karl Menger}
(1927) proved the min-cut/max-paths duality that bears his name; it is the
combinatorial heart of network-flow theory and predates the algorithmic max-flow
min-cut theorem of Ford and Fulkerson (1956) by three decades. The thread from Menger
to Ford--Fulkerson to Tarjan's linear-time low-link DFS is the throughline of
low-connectivity theory: the same facts, sharpened from existence to duality to a
linear-time certificate.
\end{historybox}

\section{Application: network reliability}
\label{ch10:reliability}

Connectivity structure is the first thing a network designer audits. Model a
communication or transportation network as a graph: vertices are routers or stations,
edges are links.

\begin{applicationbox}
\textbf{Single points of failure.} A \emph{bridge} is a link whose loss partitions the
network --- no redundancy, the most fragile element. An \emph{articulation point} is a
node whose loss partitions the network --- a router whose crash isolates a region.
Auditing a network for both is a single $\Oh(n+m)$ DFS, fast enough to run on
million-edge topologies continuously. The block-cut tree visualizes the result: each
block is a ``robust island'' that survives any single-element failure internally, and
the tree of blocks shows exactly how those islands hang together through their shared
cut vertices and bridges.
\end{applicationbox}

\begin{applicationbox}
\textbf{Reliability and augmentation.} If a network has bridges, how few links must we
add to make it $2$-edge-connected (survive any single link failure)? Contract every
$2$-edge-connected component to a point: the graph collapses to its \emph{bridge tree},
and the cheapest augmentation adds $\ceil{L/2}$ edges, where $L$ is the number of
leaves of that tree (Eswaran--Tarjan, 1976). The analogous vertex-augmentation problem
uses the block-cut tree. Both reductions depend on first computing the decomposition
this chapter builds.
\end{applicationbox}

\begin{connectionbox}
The low-link DFS is the undirected cousin of \textbf{Tarjan's strongly connected
components} algorithm (Chapter 5), which stacks vertices and uses the very same
$\mathrm{low}$ recurrence on a directed DFS tree (with cross and forward edges handled
carefully). The block-cut tree is queried with the \textbf{LCA} structures of
Chapter 23. The $k=2$ Menger statements generalize to all $k$ via the
\textbf{max-flow min-cut} theorem (Chapters 11--12), where global minimum cut and
edge/vertex connectivity become flow computations. And the ``how fragile is this
graph?'' question has a continuous, spectral answer: the second-smallest eigenvalue of
the graph Laplacian $\Lap$ (Chapters 31--32) is a quantitative measure of
connectivity, of which ``$2$-edge-connected = no bridge'' is the crude binary shadow.
Cheeger's inequality (Chapter 32) makes that link precise.
\end{connectionbox}

\section{Exercises}
\label{ch10:exercises}

\begin{exercise}
Run \textsc{Bridges} (Algorithm~\ref{ch10:alg-bridges}) by hand on the cycle $C_5$ on
vertices $1,\dots,5$, then on $C_5$ plus one chord $\{1,3\}$. Tabulate $\mathrm{disc}$
and $\mathrm{low}$ and confirm the bridge set is empty in both cases. Then delete edge
$\{4,5\}$ from the chorded graph and recompute: which edges become bridges?
\end{exercise}

\begin{exercise}
Prove that an endpoint of a bridge is an articulation point if and only if it has degree
at least $2$. Give a smallest graph where a bridge's endpoint is \emph{not} an
articulation point.
\end{exercise}

\begin{exercise}
\textbf{(Break it.)} Take the multigraph on $\{u,v\}$ with two parallel edges between
them. Show that the ``skip if $w=p$'' parent guard wrongly reports the edge $\{u,v\}$
as a bridge, and that the edge-index guard of Algorithm~\ref{ch10:alg-bridges} reports
no bridge (correctly). Explain in one sentence what cycle the second parallel edge
creates.
\end{exercise}

\begin{exercise}
Prove that the number of connected components created by deleting an articulation point
$c$ from a connected graph equals the degree of $c$ in the block-cut tree.
\end{exercise}

\begin{exercise}
\textbf{(Implement.)} Convert the recursive \textsc{BiconnectedComponents}
(Algorithm~\ref{ch10:alg-bcc}) into an iterative version using an explicit stack of
$(\text{vertex},\ \text{neighbor-iterator})$ frames, so it runs on a path of $10^6$
vertices without overflowing. Verify it emits identical blocks on the running example.
\end{exercise}

\begin{exercise}
Show that in any DFS of a connected undirected graph, the root is an articulation point
iff it has $\ge2$ children, \emph{regardless} of which vertex is chosen as root and in
what order neighbors are visited. (Hint: the \emph{set} of articulation points is
traversal-independent; only the tree shape changes.)
\end{exercise}

\begin{exercise}
\textbf{(Bridge tree.)} Describe an $\Oh(n+m)$ algorithm that contracts each
$2$-edge-connected component of $\Gph$ to a single vertex, producing the \emph{bridge
tree} whose edges are exactly the bridges of $\Gph$. Prove the result is a tree (a
forest if $\Gph$ is disconnected).
\end{exercise}

\begin{exercise}
\textbf{(Augmentation.)} Let the bridge tree of a connected graph have $L\ge2$ leaves.
Prove that adding $\ceil{L/2}$ edges suffices to make the graph $2$-edge-connected, and
that no smaller number works. (This is the Eswaran--Tarjan bound; pair leaves across the
tree.)
\end{exercise}

\begin{exercise}
\textbf{(Vertex version of low.)} Show that the recurrence used here,
$\mathrm{low}[u]=\min(\mathrm{disc}[u],\ \min_{\text{back }\{u,w\}}\mathrm{disc}[w],\
\min_{\text{child }c}\mathrm{low}[c])$, is equivalent to Definition~\ref{ch10:def-low},
and explain why using $\mathrm{low}[w]$ instead of $\mathrm{disc}[w]$ for back edges
would be \emph{wrong} for the articulation test.
\end{exercise}

\begin{exercise}
\textbf{(Extend.)} A graph is $2$-vertex-connected iff it has an open ear
decomposition. Use this to prove that every $2$-connected graph on $n\ge3$ vertices has
at least $n$ edges, with equality exactly for cycles.
\end{exercise}

\begin{exercise}
Give a connected graph in which a vertex of maximum degree is \emph{not} an articulation
point while a vertex of degree $2$ \emph{is}. Conclude that ``importance'' for
connectivity is not the same as degree centrality.
\end{exercise}

\begin{exercise}
\textbf{(The $\ge$ vs $>$ boundary.)} Construct two graphs that are identical except
for one back edge, such that the same tree edge $\{u,v\}$ satisfies
$\mathrm{low}[v]=\mathrm{disc}[u]$ in one and $\mathrm{low}[v]<\mathrm{disc}[u]$ in the
other. Verify that $u$ is an articulation point in the first but not the second, while
$\{u,v\}$ is a bridge in neither.
\end{exercise}
